% Comic Series Development Plans
% Flyxion, Independent Researcher
% Compiled: May 2026

\documentclass[12pt, a4paper]{article}

\usepackage{mathpazo}
\usepackage{geometry}
\geometry{margin=1.25in}

\usepackage{titlesec}
\usepackage{setspace}
\usepackage{hyperref}
\usepackage{microtype}
\usepackage{parskip}
\usepackage{xcolor}
\usepackage{booktabs}
\usepackage{array}
\usepackage{enumitem}

\hypersetup{
  colorlinks=true,
  linkcolor=black,
  urlcolor=black,
  pdfauthor={Flyxion},
  pdftitle={Three Comic Series Development Plans}
}

\definecolor{cosmicblue}{RGB}{20, 40, 90}
\definecolor{wormgold}{RGB}{130, 85, 10}
\definecolor{orphicpurple}{RGB}{70, 20, 90}

\titleformat{\section}{\large\bfseries\color{cosmicblue}}{}{0em}{}[\titlerule]
\titleformat{\subsection}{\normalsize\bfseries}{}{0em}{}
\titleformat{\subsubsection}{\normalsize\itshape}{}{0em}{}

\setlist[itemize]{noitemsep, topsep=4pt}

\begin{document}

\begin{titlepage}
  \centering
  \vspace*{2cm}
  {\Huge\bfseries Three Comic Series\\[0.4em] Development Plans}\\[1.5em]
  {\large\itshape An Interconnected Graphic Narrative Universe}\\[3em]
  {\large Flyxion}\\[0.5em]
  {\normalsize Independent Researcher}\\[2em]
  {\normalsize May 2026}
  \vfill
  {\small
    \textit{Series I:} Cosmic Expansion Instability (RSVP Framework)\\[0.4em]
    \textit{Series II:} Babylonian Tooth Worm Invocation\\[0.4em]
    \textit{Series III:} Orphic Mystery Comics
  }
\end{titlepage}

\newpage
\tableofcontents
\newpage

%=============================================================
\section{Overview and Thematic Universe}

Three comic series have emerged from a single sustained creative and theoretical session, each addressing the same deep question from a different civilisational vantage point: what happens when a universe that appears stable turns out to have been unstable from its very beginning, and when the creature at the end of the creation chain turns out to be the one whose hunger drives everything forward?

\textit{Cosmic Expansion Instability} translates a precise mathematical result from general relativity into a hard science fiction narrative. The flat Friedmann universe is an unstable saddle point; the cosmos we inhabit is a trajectory that has been falling away from that saddle since the first moments after the Big Bang. The series is structured around the RSVP (Relativistic Scalar-Vector Plenum) framework, which treats space not as an empty void but as a living medium --- the Lamphrodyne Plenum --- of three coupled fields: the scalar accessibility field $\Phi$, the flow field $\mathbf{V}$, and the configurational entropy field $S$.

\textit{The Babylonian Tooth Worm Invocation} adapts one of the oldest known healing texts, a Mesopotamian cuneiform incantation in which the Worm is the last creature generated by the full chain of creation from sky to mud, petitions the gods for food, is offered the roots and blood of the tooth instead of figs and apricots, and becomes the cosmological cause of dental pain. This series takes the cosmology of the incantation entirely seriously: creation is a cascade of delegated generativity, and the smallest entity at the bottom of that cascade is the one whose unsatisfied need reshapes the world.

\textit{Orphic Mystery Comics} recovers the ancient Greek cosmogonic tradition in which the universe hatches from the Primordial Egg, and the separation of fire, water, and earth into the domains of Zeus, Poseidon, and Pluto is not merely a mythological partition but a dynamical event with structural consequences for everything that follows. This series is the most cosmologically ambitious of the three, treating Orphic theogony as a genuine physics of emergence.

The three series form a thematic universe rather than a shared fictional universe. They do not share characters or settings. They share an intuition: that the interesting dynamics of any system are located not at its intended stable configuration but at the latent instabilities that were present from the moment of origin.

%=============================================================
\newpage
\section{Series I: Cosmic Expansion Instability}

\subsection{Source Material and Scientific Foundation}

The scientific basis for this series is the paper \textit{The Instability of Critical and Underdense Friedmann Spacetimes at the Big Bang as an Alternative to Dark Energy} by Alexander, Temple, and Vogler (2026, \textit{Proceedings of the Royal Society A}), synthesised with the broader RSVP framework developed by Flyxion Research. The paper establishes that the flat Friedmann spacetime ($k=0$) functions as an unstable saddle rest point $S_M$ in a self-similar dynamical system. Two eigenvalues are negative at this point --- $\lambda_1^B = -1$ (the gauge mode) and $\lambda_2^B = -\tfrac{1}{3}$ (the physical instability) --- while all eigenvalues from order $n \geq 3$ are positive (stable), following the formula $\lambda_B^n = \tfrac{2n-5}{3} > 0$. This is the Closure Depth Invariant: $d_{\text{closure}} = 2$. The universe appears infinitely complex, but all its instability content is fixed at depth two.

The canonical dynamical flow is $SM \to F \to M$: the universe departs the Friedmann saddle, traverses the family $F$ of underdense, accelerating solutions, and approaches asymptotically a Minkowski-like attractor. What we observe as dark-energy-driven acceleration is the universe falling off the saddle. There is no external push. There is only instability.

The RSVP framework enriches this picture. Space is not empty. It is the Lamphrodyne Plenum, a continuous medium of three interwoven fields. The scalar accessibility field $\Phi$ reshapes spacetime: regions of high $|\nabla\Phi|$ are metrically distant and hard to cross, while valleys of high $\Phi$ are easy corridors through which trajectories naturally flow. The accessibility metric is $ds^2_{\text{acc}} = (1 + \kappa|\nabla\Phi|^2)(dx^2 + dy^2)$. The flow field $\mathbf{V}$ carries organised bulk motion, constrained by $\Phi$ and driven by $S$. The configurational entropy field $S$ drives the system toward lower constraint. Their coupling is governed by the Lamphrodyne divergence constraint: $\nabla^{\text{acc}} \cdot \mathbf{V} = \alpha_\Phi \Phi + \alpha_S S$. Galaxies are eddies where flows of information curl into stable vortices. Expansion is the plenum relaxing from high constraint to high accessibility.

\subsection{Issue Structure}

The series is planned as eight issues, numbered to match the infographic sequence already developed.

\subsubsection{Issue 1: The Great Mistake}

The universe expands. The expansion speeds up. Scientists check the equations and find $\ddot{a} > 0$ where deceleration was expected. The response is to add a box to the universe labelled \textsc{Dark Energy} $\Lambda$. The issue closes with two voices in the dark: \textit{But what if we got it backwards? What if nothing is pushing\ldots\ but something is falling?} The tagline: \textsc{Maybe the acceleration isn't a mystery. Maybe it's an inevitability.}

\subsubsection{Issue 2: The Saddle --- Our Unstable Beginning}

The critical Friedmann spacetime is not stable ground; it is a saddle point, stable in some directions and unstable in others. The issue dramatises this through the figure of a ball perfectly balanced on a mountain pass. The tiniest nudge --- generic underdense perturbations --- sets the fall in motion. No fine-tuning required. No dark energy. Just instability. The mathematics behind the picture is displayed directly: the two negative eigenvalues, the positive spectrum for $n \geq 3$, the cosmic journey from high constraint to low constraint. The tagline: \textsc{We are not in a stable universe. We are falling. And that is why we are expanding.}

\subsubsection{Issue 3: Enter the Plenum --- Space is Alive}

The universe is not an empty void. The RSVP framework introduces the Lamphrodyne Plenum: three fields that together produce what we experience as gravity, cosmic expansion, and time. $\Phi$ shapes accessibility; $\mathbf{V}$ moves the plenum; $S$ drives relaxation. The three figures of the plenum --- the accessibility goddess, the flow god, the entropy goddess --- speak in sequence: \textit{I shape the landscape of accessibility. I carry change. I move the plenum. I drive the system toward lower constraint.} Their interaction is governed by the Lamphrodyne Functional $L$, the total energy and resistance of the medium. The tagline: \textsc{Space is not empty. Space is busy.}

\subsubsection{Issue 4: The Accessibility Landscape --- Distance Depends on the Terrain}

In RSVP, the scalar accessibility field $\Phi$ reshapes spacetime. Paths bend; distances stretch; the easiest route is not the straightest line. The issue contrasts normal Euclidean space (flat terrain, shortest path is a straight line, $ds^2 = dx^2 + dy^2$) with accessibility space (nontrivial terrain, paths curve to avoid expensive regions, $ds^2_{\text{acc}} = (1+\kappa|\nabla\Phi|^2)(dx^2+dy^2)$). Geodesics on the landscape curve naturally through valleys of high accessibility. Gravity and redshift emerge from how trajectories navigate this terrain. The tagline: \textsc{The shortest path is not always the easiest path. The universe follows accessibility.}

\subsubsection{Issue 5: The Fall of Space --- Expansion is Relaxation}

The universe does not expand into pre-existing emptiness. It relaxes, falling from high constraint to low constraint. Three epochs are depicted. In the early universe, the accessibility field $\Phi$ is highly compressed, entropy $S$ is low with extreme gradients, flow $\mathbf{V}$ is tightly constrained, and the system sits near the unstable equilibrium. In the middle universe, constraint cracks, structures emerge, voids open, and accessibility increases in patches. In the late universe, $\Phi$ smooths, barriers weaken, flows broaden, and the universe approaches maximal accessibility. What we call expansion is the plenum becoming easier to inhabit. The tagline: \textsc{Expansion is not a mystery to solve. It is a process to understand.}

\subsubsection{Issue 6: The Depth-Two Secret --- The Universe is Finite-Level Determined}

The RSVP framework reveals a profound mystery: the universe is infinitely unstable at every order, yet its full behaviour is fixed by information contained only at the second order. The eigenvalue spectrum is displayed as a table: orders 1 and 2 are unstable (negative eigenvalues), orders $n \geq 3$ are stable (positive eigenvalues, formula $\lambda_B^n = \tfrac{2n-5}{3}$). The Closure Depth Invariant $d_{\text{closure}} = 2$ means the universe is infinite in extent but finite in instability content. Beyond level 2, the tower dissolves. No further instability remains to be discovered. The tagline: \textsc{The universe appears infinitely complex. Its instability is not.}

\subsubsection{Issue 7: The Smoking Gun --- Redshift and Rotation Have a Common Source}

RSVP makes a specific observational prediction that $\Lambda$CDM cannot: redshift and polarisation rotation are correlated, because both arise from the same physical process. In RSVP, entropy gradients ($\nabla S$) drive flows ($\mathbf{V}$); when those flows develop curl ($\omega = \nabla \times \mathbf{V}$), they twist spacetime. Photons climbing the accessibility landscape gain redshift. The same vorticity rotates their polarisation. This creates a measurable correlation $C_{\beta z} \neq 0$. The $\Lambda$CDM detective treats the two clues as unrelated; the RSVP detective finds the common suspect. The test: measure polarisation rotation $\beta$ and redshift $z$ across many sources. RSVP expects a positive correlation; $\Lambda$CDM expects none. The tagline: \textsc{Follow the clues. Trust the correlation. Understand the plenum.}

\subsubsection{Issue 8: The RSVP Universe --- A Cosmos That Relaxes Within Itself}

The synthesis issue. The universe is not expanding into emptiness; it is topologically closed. Galaxies are eddies where flows of information curl into stable vortices. Entropy gradients are currents that twist spacetime and shape everything within it. Trajectories are flow lines; particles, light, and information all follow geodesics of least accessibility cost. Accessibility basins create barriers and easy paths. The fall of space is relaxation from high constraint to high accessibility. Finite depth, infinite richness: all instability lives at depth 1 and 2, infinite stable structure rests above. The tagline: \textsc{The cosmos is not expanding into emptiness. It is relaxing within itself.}

\subsection{Visual Design}

The visual grammar of the series combines hard-science infographic language with mythological illustration. Mathematical notation is treated as native visual material: eigenvalue tables, phase portraits, and accessibility metric formulae appear directly in panels. The colour palette runs from deep cosmological blue and black through gold and amber for the plenum fields. Panel borders are keyed to spacetime curvature content: flat panels for Friedmann-like regions, subtly warped borders in underdense trajectories. The three plenum fields are each given a visual personification maintained consistently across issues.

\subsection{Production Summary}

\begin{center}
\begin{tabular}{lp{9cm}}
\toprule
\textbf{Parameter} & \textbf{Specification} \\
\midrule
Format & 8 issues, prestige format, 48 pages each \\
Colour & Full colour; deep blue, gold, amber palette \\
Scientific basis & Alexander, Temple \& Vogler (2026, \textit{Proc.\ R.\ Soc.\ A}) + RSVP/CLIO/TARTAN framework \\
Panel geometry & Curvature-responsive borders \\
Visual register & Hard science infographic + mythological illustration \\
Publisher label & Flyxion Research \\
\bottomrule
\end{tabular}
\end{center}

%=============================================================
\newpage
\section{Series II: The Babylonian Tooth Worm Invocation}

\subsection{Source Material}

The source is one of the oldest surviving healing texts: the Babylonian incantation against the tooth worm, preserved in cuneiform and dating to at least the second millennium BCE. Its structure is a cosmogonic cascade. When the gods created the heavens, the heavens created the earth, the earth created the rivers, the rivers created the canals, the canals created the marshes, and the marshes created the Worm. At the bottom of every chain of creation sits an entity that was created last and given nothing to eat. The Worm petitions Anu and Ea. Anu offers the ripe fig and the apricot. The Worm refuses: \textit{Place me and let me dwell between tooth and gum --- I will suck the tooth's blood and devour the gum.} Ea's response, the formula of expulsion, becomes the healer's tool.

The concept art has already developed this into a full eight-issue series under the banner \textit{Babylonian Tooth Worm Invocation}, with two distinct cover art iterations demonstrating the visual range of the series. A complete interior spread for Issue 5 (\textit{The House of Water}) has also been developed.

\subsection{Issue Structure}

\subsubsection{Issue 1: The Birthchain of Creation}

The cosmogonic cascade is rendered in full. From sky to mud, a chain of creation unfolds, each level generating the next, each delegation diffusing responsibility downward. At the end of the chain, the Last Child awakens. The tagline: \textit{From sky to mud --- a chain unfolds. At the end, the Last Child awakens.}

\subsubsection{Issue 2: The Worm's Complaint / The Complaint Before the Sun}

The Worm approaches the divine assembly. The complaint is cosmologically structured: \textit{What have you given me to eat? What have you given me to suck?} The response from Anu is the ripe fig and the apricot. But the Last Child desires more than fruit. Two cover versions have been developed: one staging the petition in a neoclassical divine court, the other using an older, more archaic visual grammar.

\subsubsection{Issue 3: Rebellion in the Mouth of the World / Between Tooth and Gum}

The Worm speaks its counter-demand. Between tooth and gum --- the smallest space in the world --- becomes a kingdom. The order of things is threatened. The tagline: \textit{So speaks the Last Child. And the order of things is threatened.} / \textit{Between tooth and gum --- the smallest space becomes a kingdom.}

\subsubsection{Issue 4: The Accessibility War}

The divine enforcement apparatus mobilises. The smallest crack can unravel the greatest order. Tiny spaces, great structures, only one can hold. This issue draws the explicit thematic connection to the RSVP accessibility landscape: regions of the cosmic plenum where constraint is violated produce instability, just as a single Worm in the wrong cavity destabilises the entire structure.

\subsubsection{Issue 5: The House of Water}

The most fully developed issue in the current concept art. Below the luminous realms of sky and star, below the depths and rivers, lies a domain older than kings and older than memory --- the House of Water. Here the laws of balance are kept, the currents of creation are measured, and the boundaries of possibility are guarded. The Keepers of Flow and Form, the memory of what is, the watchers of what may be, receive the report: there is one who has broken the order. Born of the mud, the last of the chain, it cries for more than was given. It seeks the small cracks between the solid things. And where it goes, pain follows. The House acts. The Worm will be driven out. The tagline: \textit{So it is written. To be continued.}

\subsubsection{Issue 6: The Tooth of Heaven}

Beyond gods and stars stands the structure that holds them all: a crystal that holds the cosmos together, and a crack that could break it apart.

\subsubsection{Issue 7: The Trial of the Worm / The Trial of the Last Child}

The Worm is brought before the divine tribunal. \textit{Why did you do this?} The answer given in the concept art is the most resonant line of the series: \textit{Because there was nowhere left for me to live.} The question the issue poses: is this a crime or a cry of exclusion? Rebel or child? The truth is older than their law.

\subsubsection{Issue 8: The Secret of the Last Child}

The synthesis issue. Without disruption, there is no change. Without the Last Child, there is no universe. The Worm's action, its refusal of the assigned food and its demand for the blood of the tooth, was not rebellion but necessity. The secret is that the universe requires its most disruptive element. The tagline: \textit{Without disruption, there is no becoming. The Last Child teaches the universe to move.}

\subsection{Visual Design}

Two visual registers have been developed in the concept art and both should be maintained as variant covers or as a visual evolution across the series. The first is a classical adventure comic aesthetic with lush colour and dynamic figure composition. The second is a more archaic visual grammar inspired by ancient Near Eastern conventions: strict bilateral registers, flat profile figures, hierarchical scaling, ground lines as panel dividers. The interior spread for Issue 5 demonstrates a successful synthesis: the House of Water is rendered with the grandeur of a Mesopotamian palace complex while the panel structure follows contemporary adventure comics conventions.

The completed interior spread for Issue 5 is the most valuable design asset currently available and should serve as the visual bible for the remainder of the series.

\subsection{Production Summary}

\begin{center}
\begin{tabular}{lp{9cm}}
\toprule
\textbf{Parameter} & \textbf{Specification} \\
\midrule
Format & 8 issues, standard format, 32 pages each \\
Colour & Full colour; deep ochre, lapis, terracotta, gold \\
Source text & Babylonian Tooth Worm Incantation (c.\ 2nd millennium BCE) \\
Issue dates (in-world) & August 1982 -- March 1983 \\
Cover price (in-world) & \$1.25 US / \$1.90 CAN \\
Publisher label (in-world) & Babylonian Tooth Worm Invocation \\
Visual basis & Ancient Near Eastern art + classical adventure comics \\
\bottomrule
\end{tabular}
\end{center}

%=============================================================
\newpage
\section{Series III: Orphic Mystery Comics}

\subsection{Source Material}

The ancient Greek Orphic theogony is one of the most cosmologically ambitious mythological systems in the Western tradition. In its most developed form, the universe begins in an absolute void from which the Primordial Egg emerges, containing within it the potential for all differentiation. From the Egg hatches Phanes, the first light, the first form, the principle of manifestation. From Phanes proceeds the ordered cosmos: fire ascends to become the domain of Zeus, the waters flow to become the domain of Poseidon, and the heavy descends to become the domain of Pluto. This is not merely a partition of the cosmos among three brothers; it is a dynamical account of how differentiation itself happens, how a single undifferentiated origin produces the structured multiplicity we inhabit.

The concept art has established the visual and narrative register of this series through a single issue cover: \textit{Orphic Mystery} Issue 3, published under the imprint \textit{Ancient Wisdom Comics} (approved by the Oracles Code Authority, \$1.50 US / \$1.90 CAN / UK \pounds1.10). The tagline: \textit{Before gods. Before worlds. Before memory. There was the Egg.} / \textit{From chaos came form! From the Egg came the World!}

\subsection{Concept and Tone}

This series occupies a different register from the other two. Where the RSVP series is hard science fiction and the Tooth Worm series is mythological horror, Orphic Mystery is cosmogonic epic: vast in scale, reverent in tone, and formally reminiscent of the Silver Age comics that established the superhero cosmology genre (Kirby's New Gods, O'Neil's Green Lantern). The three domains --- fire/Zeus, water/Poseidon, earth/Pluto --- are treated as dynamical attractors rather than merely spatial territories, and the series is interested in what happens at the boundaries between them.

\subsection{Issue Structure}

The series is planned as an ongoing title under the \textit{Ancient Wisdom Comics} imprint. The single developed issue cover (Issue 3, February) suggests a series already in progress, with the first two issues covering the pre-Egg void and the hatching of Phanes. The known issues are as follows.

\subsubsection{Issues 1--2: Before the Egg / The Hatching}

The absolute void precedes even the Egg. Issues 1 and 2 establish the cosmogonic conditions: the Orphic Night, the first stirring, and the hatching of Phanes. These issues would be the most formally experimental of the series, operating without conventional panel structure in some sequences to convey the absence of differentiated spacetime.

\subsubsection{Issue 3: The Birth of the Cosmos}

The confirmed issue. Fire ascends. Waters flow. Earth descends. Order is born. Zeus, Poseidon, and Pluto each claim their domain in a three-way differentiation event that is simultaneously a physical phase transition and a political act. The cover depicts the moment of maximum clarity: Phanes radiant at the centre, Zeus as flame on the left, Poseidon as ocean on the right, Pluto as stone below. The tagline: \textit{From chaos came form! From the Egg came the World!}

\subsubsection{Issues 4 onwards: The Orphic Mysteries Proper}

The later issues of the series would concern the ongoing dynamics of a differentiated cosmos: what happens when fire, water, and earth interact at their boundaries, what the Orphic initiatory tradition understood about those boundary dynamics, and how the human soul navigates a cosmos structured by the three-domain differentiation. The Orphic gold tablets, the Orphic Hymns, and the Rhapsodic Theogony provide narrative material for an extended run.

\subsection{Visual Design}

The confirmed issue cover establishes a visual language that is simultaneously indebted to Silver Age comics aesthetics and to ancient Greek vase painting conventions. The use of strict figural hierarchy (Phanes larger and more luminous than the three gods, the three gods in distinct colour registers: gold/flame for Zeus, blue/wave for Poseidon, purple/stone for Pluto), the dramatic border design with the starburst callout, and the physical cover format (barcode, approval seal, price box) all locate the series firmly within the visual tradition of genre comics while treating mythological subject matter with genuine seriousness.

The \textit{Ancient Wisdom Comics} imprint should be developed as a distinct brand identity from the Flyxion Research imprint used for the RSVP series and the Babylonian Tooth Worm imprint used for Series II.

\subsection{Production Summary}

\begin{center}
\begin{tabular}{lp{9cm}}
\toprule
\textbf{Parameter} & \textbf{Specification} \\
\midrule
Format & Ongoing monthly title \\
Colour & Full colour; gold, flame, ocean blue, deep purple \\
Source material & Orphic Theogony, Orphic Hymns, Rhapsodic Theogony \\
Issue confirmed & Issue 3 (February, in-world) \\
Cover price & \$1.50 US / \$1.90 CAN / UK \pounds1.10 \\
Publisher label & Ancient Wisdom Comics \\
Visual basis & Silver Age comics + ancient Greek figural conventions \\
\bottomrule
\end{tabular}
\end{center}

%=============================================================
\newpage
\section{Cross-Series Thematic Architecture}

\subsection{The Three Creation Chains}

Each series is built around a creation chain in which the generative cascade produces, at its final step, an entity whose existence destabilises everything above it. In the RSVP series, the cascade of physical fields produces a universe whose instability at depth two propagates consequences through all higher levels. In the Tooth Worm series, the cascade from sky to marsh produces a creature whose hunger the cosmic order cannot satisfy and whose demand for territory reorganises the experience of embodied creatures. In the Orphic series, the cascade from the primordial void through the Egg to Phanes to the three-domain differentiation produces a structured cosmos in which the boundaries between domains are the sites of maximum dynamical interest.

The three chains share a formal structure: a generative sequence that was supposed to produce a stable terminus, a terminal entity or configuration whose properties were not anticipated by the levels above it, and a reorganisation of the entire system in response.

\subsection{The Three Depth-Two Moments}

The RSVP framework gives a precise name to the critical level of a generative hierarchy: the Closure Depth Invariant $d_{\text{closure}}$. For the universe, this is 2. The instability that determines everything lives at level 2, not at higher levels. This concept can be read across all three series. In the Tooth Worm series, the Worm is the second-order creature: the marshes created it (one level above), and the marshes were created by the canals (two levels above). Its instability is a depth-two phenomenon in exactly the RSVP sense. In the Orphic series, the three-domain differentiation is the depth-two event: the void created the Egg (one level), and the Egg created Phanes (two levels), and from Phanes the entire structure of the cosmos derives. The differentiation at depth two determines all subsequent structure.

\subsection{The Common Intuition}

The three series share a single intuition that can be stated briefly: the interesting dynamics of any sufficiently complex system are not located at its intended stable configuration but at the latent instability that was present from the moment of origin and that drives every subsequent development. The universe we live in is not the intended universe; it is the one that fell off the saddle. The pain we experience in our teeth is not a malfunction of the body; it is the cosmological consequence of the Worm's placement. The structure of the cosmos is not the void that preceded it; it is what emerged from the Egg's crack.

In each case, the origin story is also an instability story. And in each case, what looks like a mistake or a failure or a disruption turns out to be the condition of everything actual.

%=============================================================
\newpage
\section{Immediate Development Priorities}

\subsubsection{Series I: Cosmic Expansion Instability}

Issues 1 through 8 have concept art infographics completed. The priority is scripting Issue 1 in full, establishing the voice and panel rhythm of the series. The phase portrait spread planned for Issue 2 should be designed next, as it will set the visual conventions for all mathematical content across the run. The three plenum field personifications ($\Phi$, $\mathbf{V}$, $S$) introduced in Issue 3 need character design sheets establishing their visual consistency.

\subsubsection{Series II: The Babylonian Tooth Worm Invocation}

The Issue 5 interior spread (\textit{The House of Water}) is the most complete interior asset available and should serve as the visual bible. The two cover design iterations provide a range of visual registers to choose between for the final series look. Priority is scripting Issues 1 and 2, which establish the cosmogonic register before the Worm speaks. The Worm character design requires specific attention: it must be small, precise, and visually legible as a creature with genuine cosmological standing.

\subsubsection{Series III: Orphic Mystery Comics}

Only Issue 3 has confirmed concept art. The immediate priority is developing issues 1 and 2 retroactively, establishing the pre-Egg and hatching sequences. The \textit{Ancient Wisdom Comics} imprint identity --- logo, design conventions, the Oracles Code Authority seal --- should be formalised as a style guide. The visual logic of the three-domain colour register (gold/blue/purple for Zeus/Poseidon/Pluto) is established and should be maintained consistently.

\subsubsection{Cross-Series}

The thematic connection between the Closure Depth Invariant and the depth-two structure of the other two series should be developed into a one-page design element suitable for inclusion as supplementary material in each series. This cross-series document also serves that purpose and should be updated as the series develop.

%=============================================================
\newpage
\section{Critical Reading and Refined Architecture}

\subsection{The Boundary Motif: A Deeper Common Structure}

A critical reading of the three series in parallel reveals a structural pattern more fundamental than the creation-chain and instability motifs already identified. All three series are ultimately stories about boundaries: what defines them, what happens when they are violated or negotiated, and why the violation is not a failure of the system but the condition of everything that follows.

In the RSVP series, accessibility gradients are literally the boundaries of configuration space. Regions where $|\nabla\Phi|$ is large are metrically distant, hard to cross, effectively forbidden. The plot of each issue is the story of a trajectory navigating --- or failing to navigate --- those boundaries. The Friedmann saddle is an unstable boundary between overdense collapse and underdense acceleration. The Depth-Two Secret is the discovery that only two levels of boundary structure determine everything that follows at higher levels. The Smoking Gun is the detection of a boundary being crossed: entropy gradients curling into vorticity, spacetime itself being twisted by the transit.

In the Tooth Worm series, the critical boundary is announced explicitly and spatially: the gap between tooth and gum. The Worm's cosmological complaint is a boundary negotiation. It was created at the bottom of a generative cascade and offered a food source that belongs to a different part of the creation order. Its counter-demand --- to dwell precisely at the boundary between tooth and gum, to inhabit the crack rather than the substance --- is the move that defines the series. The House of Water is the cosmic institution whose function is to maintain the boundaries of what may be. The trial of the Worm is not a criminal proceeding but a boundary adjudication.

In the Orphic series, the entire cosmogony is a sequence of boundary-creation events. The void is the condition before any boundary exists. The Egg is the first boundary: an inside and an outside. The hatching of Phanes is the first boundary-crossing. The differentiation into the domains of Zeus, Poseidon, and Pluto is the creation of the three fundamental boundaries that structure all subsequent reality. The Orphic initiatory tradition, which provides narrative material for the later issues, is explicitly concerned with the navigation of boundaries between mortal and divine, embodied and disembodied, remembered and forgotten.

This convergence on the boundary motif is not accidental. Instability, by its mathematical definition, is the condition in which a small displacement across a boundary in configuration space produces a qualitatively different trajectory. The Worm found a boundary no one had mapped. The universe crossed the boundary of the saddle and has been travelling away from it ever since. The Orphic cosmos is the result of the primordial boundary's first instantiation. In all three cases, the story begins when something crosses a boundary it was not expected to cross.

\subsection{Reframing: Three Mythologies of Accessibility}

This convergence suggests a reframing of the entire publishing project. Rather than three comic series that share a thematic architecture, what is being developed is a single intellectual project expressed in three historical and cosmological registers: three mythologies of accessibility.

The word accessibility is precise here and connects deliberately to the RSVP framework. An accessibility mythology is a narrative that takes seriously the question of which regions of a possibility-space are reachable, which are forbidden, and what happens when a trajectory discovers a passage that was presumed closed. The RSVP series poses this question in the language of twenty-first century mathematical physics. The Tooth Worm series poses it in the language of second-millennium Babylonian cosmology. The Orphic series poses it in the language of ancient Greek theogony. Each gives a different answer to the same structural question, and the differences between the answers are as instructive as the similarities.

In the RSVP framework, accessibility is a continuous field: $\Phi$ varies smoothly, barriers are gradients, and forbidden regions are merely expensive. In the Babylonian framework, accessibility is negotiated: the Worm petitions, the gods respond, a boundary is assigned by divine decree rather than by geometry. In the Orphic framework, accessibility is creative: the differentiation of domains does not merely describe which regions are reachable but actively creates the topology upon which future trajectories become possible. The void before the Egg has no accessibility structure because it has no structure at all; the cosmos after the differentiation is a fully articulated manifold of possibility.

These three epistemologies of accessibility --- continuous/geometric, negotiated/institutional, and creative/topological --- are not reducible to one another. That is precisely what makes the three series worth doing separately rather than collapsing into a single franchise.

\subsection{Series Strengths and Remaining Work}

The RSVP series is the most structurally complete. Its eight issues constitute a genuine narrative arc from anomaly through discovery through synthesis, and the mathematical backbone provides an unusually stable foundation for long-form worldbuilding. The phase portrait of the STV-ODE is not merely an illustration; it is a map of the story's possibility space, and every issue is a different trajectory through that map. The primary remaining work is scripting and the character design for the three plenum personifications.

The Tooth Worm series has the strongest emotional core. The line \textit{Because there was nowhere left for me to live} is the emotional spine that could carry the entire project. The Worm's complaint is immediately legible regardless of the reader's familiarity with Mesopotamian cosmology or dynamical systems theory. The primary remaining work is developing the Worm's visual design and scripting the opening two issues, which must establish the cosmogonic register with enough grandeur to make the Worm's eventual appearance feel like an arrival rather than an intrusion.

The Orphic series is presently the least developed in terms of narrative architecture, though it may have the largest visual potential. The transition from undifferentiated void to differentiated cosmos naturally produces the kind of large-scale imagery that a cosmogonic epic requires. What the series currently lacks is a character-driven through-line equivalent to the Worm's complaint or the RSVP surveyor's instruments. The Orphic initiatory tradition offers candidate figures: the soul navigating the underworld, the initiate learning the password that grants access to the waters of memory rather than the waters of forgetting. One of these figures, or a newly invented one occupying a similar structural position, would give the later issues of the series the same kind of ground-level perspective that the Tooth Worm provides in Series II and the surveyor provides in Series I.

\subsection{Imprint Architecture}

The three series warrant three distinct imprint identities, already partly established in the concept art. \textit{Flyxion Research} for the RSVP series positions it as a scientific publishing event as much as a comics event, consistent with the series' debt to Alexander, Temple, and Vogler and to the RSVP theoretical framework. \textit{Babylonian Tooth Worm Invocation} as imprint name for Series II is unusual but defensible: it names the source text directly and signals to the reader that the series takes its mythology seriously rather than using it decoratively. \textit{Ancient Wisdom Comics} for Series III, with the \textit{Oracles Code Authority} seal as a parody-homage to the Comics Code Authority, establishes a Silver Age visual register that is both playful and genuinely reverent toward its source material.

The question of whether to present these three imprints as facets of a single publishing house or as entirely independent lines is a marketing and positioning decision rather than a creative one. The argument for unification is discoverability: readers who find one series can be directed to the others. The argument for separation is tonal integrity: readers who come to \textit{Babylonian Tooth Worm Invocation} for mythological horror may not want to be immediately redirected to a hard-SF cosmology series. The recommended approach is soft unification --- a common design element, perhaps the Closure Depth diagram, that appears as a hidden or backmatter element in all three series, visible to readers who look but not foregrounded in the main experience of any individual title.

\subsection{The Spine}

If the three series are mythologies of accessibility, then the spine of the entire project can be stated in a single sentence that does not depend on knowing any of the mathematics, mythology, or cosmogony involved:

\begin{quote}
\textit{Every system has a region it believes nothing can reach. The most important moments in any history are when something does.}
\end{quote}

This sentence works for the Friedmann saddle, for the gap between tooth and gum, and for the crack in the Primordial Egg. It is the publishing spine, the cross-series tagline, and the thematic answer to the question all three series are asking.

%=============================================================
\newpage
\section{Editorial Notes: Structural Interventions}

\subsection{The Orphic Initiate}

The Orphic series currently operates entirely at the scale of fundamental forces: Phanes, Zeus, Poseidon, Pluto. That is cosmogony from the perspective of the forces themselves, which sustains an issue or two but cannot carry an ongoing. The fix is already present in the source material. The Orphic gold tablets --- small laminated sheets placed in the graves of initiates, containing passwords for navigating the underworld --- are explicitly a reader-surrogate document. They tell a soul which roads to take, which guardians to address, which water to choose. The soul that must navigate the newly differentiated cosmos, choosing the waters of Memory over the waters of Forgetfulness, learning the passwords that grant passage across the domain boundaries, is exactly the ground-level protagonist the series requires.

This figure need not be named or fixed biographically. The Orphic tradition already treats the soul navigating the underworld as a generic figure, a vessel for the reader's own navigation. A rotating ``soul'' figure across issues, or a single initiate whose identity is gradually revealed, would give the later issues the same character-driven stakes that the Worm provides in Series~II and the surveyor provides in Series~I, without requiring the cosmogonic scale of the opening issues to be reduced. The boundary crossings of the initiate become the reader's entry point into a cosmos whose fundamental boundaries were established in Issues~1 through~3.

The practical consequence for scripting: Issues~1 and~2 (pre-Egg void and hatching) can be written at cosmogonic scale, since no ground-level protagonist yet has anywhere to be. Issue~3 (confirmed cover, the three-domain differentiation) introduces the domains. Issue~4 is where the initiate first appears, inheriting a cosmos already fully structured, and must learn to read its boundaries. This sequence mirrors the structure of the Orphic initiation itself: the cosmogony is recited to the initiate precisely so that they understand the topology they are navigating.

\subsection{The RSVP Detective}

The three plenum field personifications ($\Phi$, $\mathbf{V}$, $S$) carry the philosophical content of the RSVP series but do not provide a reader surrogate capable of discovery, error, or consequence. Issue~7's framing of the RSVP detective versus the $\Lambda$CDM detective is the seed of a solution. That figure --- a theoretical cosmologist navigating the accessibility landscape as both external object of study and internal mental geography --- should be developed as a recurring protagonist across the series.

The structural opportunity is precise. The accessibility landscape of the RSVP framework is not merely a physical description; it is also a cognitive one. A cosmologist whose mental model of the universe is itself a kind of accessibility landscape, who experiences the Friedmann saddle as a kind of intellectual equilibrium point she has been trained to inhabit and must learn to fall away from, gives the series an internal journey that runs parallel to the external cosmological one. Her encounter with the smoking-gun correlation in Issue~7 becomes not merely a scientific discovery but a personal reckoning: the universe she was taught to inhabit is not the universe she lives in.

This character also provides the series with dialogue. The field personifications speak in the register of physical law; the detective speaks in the register of a person trying to understand physical law. The tension between those two registers --- between the universe as it is and the universe as it has been modelled --- is the dramatic engine the series currently lacks at the human scale.

\subsection{The Silent Heraldic Page}

The soft-unification strategy outlined in Section~7.4 should be made concrete and visual rather than remaining a design recommendation. Each issue of every series across all three imprints should carry a single silent back-page element: a wordless one-page comic consisting of three parallel panels, depicting the same structural event at three different scales and in three different visual registers.

Panel one, in the visual language of the RSVP series: the Friedmann saddle, a ball at the mountain pass, the first deviation beginning.

Panel two, in the visual language of the Tooth Worm series: the gap between tooth and gum, the worm finding it, the first entry.

Panel three, in the visual language of the Orphic series: the Primordial Egg, the first crack, the first light.

No words. No labels. The three images speak the same event in three languages simultaneously. Over the course of the publishing run, this page changes subtly --- the ball has moved slightly further down the slope, the worm has advanced slightly further into the gap, the crack in the egg has opened slightly wider --- until in the final issues of each series, the panel corresponding to that series shows the terminus: the universe at the Minkowski shore, the worm condemned and assigned, Phanes fully emerged. The other two panels continue their own trajectories, pointing toward their own conclusions.

This page functions as a heraldic emblem. A reader who picks up any single issue from any of the three series receives the cross-series intuition without requiring the other two series. A reader who collects all three series watches three simultaneous progressions. The Closure Depth diagram identified in Section~6 should be incorporated into the corner of this page as a small recurring mark, analogous to a publisher's device.

\subsection{Production Sequence}

Given the existing asset distribution --- eight completed infographic spreads for Series~I, the House of Water interior spread and two cover iterations for Series~II, and a single confirmed cover for Series~III --- the recommended scripting sequence is as follows.

Script Series~I Issue~1 first. The RSVP series establishes the scientific register that provides the intellectual permission for the other two series to be as cosmologically ambitious as they are. A reader who understands that a respectable mathematical physics paper establishes the saddle instability as a genuine result enters the Tooth Worm and Orphic series with a different cognitive preparation than one who does not. Series~I Issue~1 is also the point of maximum accessibility for a general reader, since it opens with the anomaly --- the universe is accelerating and we do not know why --- rather than with the solution.

Script Series~II Issue~1 and the Worm character design second. The emotional spine of the entire project is \textit{Because there was nowhere left for me to live}, and that line lands in Issue~7. For it to land, the Worm must be established as a credible cosmological agent from Issue~1. The character design work is therefore not decorative but structural: the Worm's visual register determines whether the reader treats it as a creature with genuine standing in the creation order or as an illustrative device.

Retro-build Series~III Issues~1 and~2 after the initiate figure is designed, not before. The cosmogonic scale of those issues is best determined once the ground-level perspective of the initiate is established, because the initiate's eventual entry point in Issue~4 sets the ceiling on how abstract the opening issues can be while remaining readable.

\subsection{The Preferred Reading Order}

The document has treated the three series as parallel and independent, suitable for discovery in any sequence. A different logic is also available, which should be acknowledged here even if it is not enforced in the publication strategy.

The three series can be read as a single argument in three movements, ordered not chronologically but conceptually. \textit{Orphic Mystery Comics} establishes the origin of boundaries: before the Egg, no structure; after the crack, a fully articulated topology. \textit{Cosmic Expansion Instability} establishes the dynamics within boundaries: how trajectories navigate an accessibility landscape, how the saddle is unstable, how a universe falls. \textit{The Babylonian Tooth Worm Invocation} establishes what inhabits the unmarked boundary: the gap that no one assigned, the creature that no one provisioned, the instability that lives permanently at the edge of the structure rather than within it.

Read in this order --- Orphic, then RSVP, then Tooth Worm --- the three series constitute a complete cosmological argument: boundaries are created, dynamics occur within them, and the most consequential entity is the one that finds the boundary no one thought to map. A collected edition or a box set would benefit from this sequencing made explicit, with a short connecting essay that articulates the argument the three series are collectively making.

\vspace{2em}
\noindent\rule{\linewidth}{0.4pt}

\vspace{1em}
\noindent\textit{Document prepared by Flyxion, Independent Researcher. May 2026. Based on creative development sessions reviewing \textit{The Instability of Critical and Underdense Friedmann Spacetimes at the Big Bang as an Alternative to Dark Energy} (Alexander, Temple \& Vogler, 2026), the Babylonian Tooth Worm Incantation, and the Orphic theogonic tradition, with reference to concept art produced in the Cosmic Expansion Instability creative session.}

\end{document}
